% !PRINTED 75
\weberpage{75}
%<<CONT
is. To determine the conditions for this case, we extend somewhat
further the concept of the \gesperrt{resolvents}, as we have already
considered it in \S.\,17.

We let $n$ run through the sequence of the numbers of the group $\mathfrak{N}$,
and denote by $\chi(n)$ one of the characters of this
group. The extended resolvents are then, if $r$ is an
$m$th root of unity:
\weq{5}{(\chi, r) = \sum^{n} \chi(n) r^n.}
If we understand by $\mathfrak{B}$ the group reciprocal to $\mathfrak{A}$ and let
$b$ run through the numbers of $\mathfrak{B}$, then by \S.\,12, 7. there is
a number of characters $\chi_b$, equal to the degree of $\mathfrak{B}$, which are
distinguished by the fact that
\[
\chi_b(a) = 1
\]
for every number $a$ from the group $\mathfrak{A}$. From this it then follows,
by the definition of the characters \S.\,11, (4), for every $n$:
\weq{6}{\chi_b(a n) = \chi_b(n),}
that is, the character $\chi_b$ has one and the same value for all
numbers of each subsidiary group $\mathfrak{A} n_1, \mathfrak{A} n_2 \ldots$ Accordingly
the resolvent $(\chi_b, r)$ will depend only on the periods $\eta$, and
will take the expression:
\weq{7}{(\chi_b, r) = \eta + \chi_b(n_1) \eta_1 + \cdots + \chi_b(n_{e-1}) \eta_{e-1}.}
If $\eta$ belongs to a group $\mathfrak{A}'$, then the conjugate
numbers $\eta_1, \eta_2, \ldots \eta_{e-1}$ belong to the same group (because $\mathfrak{A}'$ is a
normal divisor of $\mathfrak{N}$, Bd.\ I, \S.\,154). If therefore $a'$ is any
number from $\mathfrak{A}'$, the quantities $\eta, \eta_1, \eta_2, \ldots \eta_{e-1}$ remain
unchanged by the substitution $(r, r^{a'})$, and by (7) we also have
\weq{8}{(\chi_b, r) = (\chi_b, r^{a'}).}
On the other hand, one obtains from (5), if one bears in mind that
$a' n$ runs through the whole group $\mathfrak{N}$ together with $n$,
\[
\begin{aligned}
(\chi_b, r) &= \sum^{n} \chi_b(n a') r^{n a'} = \chi_b(a') \sum^{n} \chi_b(n) r^{n a'}\\
&= \chi_b(a') (\chi_b, r^{a'}),
\end{aligned}
\]
hence by (8):
\weq{9}{(\chi_b, r) = \chi_b(a') (\chi_b, r).}
If $\mathfrak{A}$ is not identical with $\mathfrak{A}'$, but is a proper divisor
of $\mathfrak{A}'$, then the group $\mathfrak{B}'$ reciprocal to $\mathfrak{A}'$ is a proper divisor
of $\mathfrak{B}$ (by \S.\,12, 8.); if therefore $b$ is a number in $\mathfrak{B}$ which does not
%CONT>>
